\section*{Ethics Statement}

This research predicts antibiotic resistance in \textit{Escherichia coli} directly from bacterial
genomic data, and no human participants, patient records, animal subjects or new laboratory
experiments were involved at any stage. All genomic and phenotypic data used were bacterial
rather than human, and every dataset was obtained from a public repository released for open
reuse by the research group that produced it~\cite{morin2024ecoli, york2025amr, Ren2022}. No new
sequencing, culturing or wet-lab work was carried out, so every result reported here is computed
from already-published, de-identified strain-level records. No patient identifiers, clinical
notes, demographic information or treatment records were accessed, requested or processed at any
point; the only strain-level fields read by the pipeline were the genome name, sequencing
accession, variant sample identifier, phylogroup assignment and antibiotic susceptibility label.
Because the project involved no living subjects and no data collection beyond downloading public
bacterial datasets, it required no animal or human trials and no institutional ethics board
approval.

Two further commitments shape how the work is reported. The system developed here is a research
prototype that has not undergone clinical validation and is not fit for deployment in a
diagnostic or treatment setting, and this is stated explicitly wherever the clinical relevance
of the work is discussed, so that no result is presented as ready for clinical use. Where
findings, code or reasoning from other published work were used they are cited in full, and the
two dataset releases this research draws on are credited separately rather than treated as a
single source. A fuller discussion of the societal, environmental and ethical implications of
this work, covering data provenance, fairness across bacterial populations, dual-use risk and
research integrity practices, is given in Chapter~\ref{ch:req} from
Section~\ref{sec:societal} onward.

\pagebreak
